\section{Tight Lower Bounds using $\Pi$-Matroids}
\label{sec:tight_absolute}
In this section we use \Cref{thm:emb} and \Cref{cor:deter} to prove \Cref{thm:main}. We rely on the following reduction.

\begin{definition}
	Given an \textnormal{EMB} instance $I =(E,\cI,c,k,T)$ and non-trivial  \textnormal{MOL} problem $P \in \cP \setminus \{(\min,\IS,\leq)\}$ where $P = (\opt,\cF,\triangleleft)$, define $R_P(I) = (E,\cI,v_I,w_{I,P},B_{I,P})$ as the {\em reduced} $P$ instance of $I$ as follows.
	\begin{enumerate}
		\item Define the following auxiliary variable \begin{align*}
		%	n & : \text{The variable representing a natural number.} \\
			d(P) = & \begin{cases}
				0 & ~~\text{if  } \left(\opt = \max, \triangleleft = \leq \right) \text{ or } \left(\opt = \min, \triangleleft = \geq \right) \\
			1 & ~~\text{otherwise.}
			\end{cases}
		\end{align*}
		
		\item For all $e \in E$ define $v_I(e) = 2 \cdot c(E)+c(e)$. 
		\item For all $e \in E$ define $w_{I,P}(e) = 2 \cdot c(E)+c(e) \cdot (-1)^{d(P)}$.
		\item Define $B_{I,P} = k \cdot 2 \cdot c(E)+T \cdot  (-1)^{d(P)}$. 
	\end{enumerate}
\end{definition}



%, our lower bounds for {\sc budgeted matroid independent set}, {\sc exact weight matroid independent set}, {\sc constrained minimum weight basis of a matroid}, and {\sc knapsack cover with a matroid}. All proofs use reductions from the $\Pi$-matroid problem.

%In all of the constructions, we are able to use (with slight modifications) the same family of {\em standard weight} functions and {\em standard budgets}. 


%Specifically, for every $n \in \mathbb{N}_{>0}$ define the standard weight function of $n$ as $w_{n}:[n] \rightarrow \mathbb{N}$ such that for all $i \in [n]$
% \begin{equation}
%	\label{eq:w}
%	w_{n}(i) = i+2 \cdot \textnormal{\textsf{sum}}([n]). 
%\end{equation}
%Also, define the standard budget of  $n,k,\alpha \in \mathbb{N}_{>0}$  as 
%\begin{equation}
%	\label{eq:B}
%	B_{n,k,\alpha} = \alpha+2 k \cdot \textnormal{\textsf{sum}}([n]). 
%\end{equation}

%In deriving the lower bounds for our problems (in
%\Cref{sec:abs_BM,sec:abs_EMI,sec:abs_CMB,sec:abs_KCM}), we use some general results, as stated in \Cref{lem:OPTu,lemma:if,lemma:onlyif}.
The following lemma gives a polynomial bound on the total weight of any subset of the ground set of a reduced instance. 

\begin{lemma}
	\label{lem:OPTu}
%	For every $n \in \mathbb{N}_{>0}$ and $S \subseteq [n]$, it holds that $w_{n}(S) \in \{0,1,\ldots, 3 \cdot n^3\}$. 
	Given an \textnormal{EMB} instance $I =(E,\cI,c,k,T)$ and non-trivial  \textnormal{MOL} problem $P \in \cP \setminus \{(\min,\IS,\leq)\}$ where $P = (\opt,\cF,\triangleleft)$, for every $S \subseteq E$ it holds that $v_{I}(S) \in \{0,1,\ldots, 2 \cdot c(E) \cdot k\}$.  
\end{lemma}

\begin{proof}
	By \eqref{eq:w} we have that $w_{n}(S) \in \mathbb{N}$. Furthermore,  $0 = |\emptyset| \leq |S| \leq \textsf{sum}(S) \leq w_{n}(S)$, and  $$w_{n}(S) \leq n \cdot \max_{i \in [n]} w_{n}(i) = n \cdot \max_{i \in [n]} i+ n \cdot (2 \cdot \textsf{sum}([n])) \leq n^2+ 2 \cdot n^3 \leq 3 \cdot n^3.$$ 
\end{proof}


We use the sets $\mathcal{J}_{n,k}$, $\mathcal{K}_{n,k,\alpha}$, $\mathcal{L}_{n,k,\alpha}(\Pi)$ and $\cI_{n,k,\alpha}(\Pi)$  as in \Cref{def:Matroid}. As in \Cref{sec:Pi} we often omit the subscripts $n,k$ and $n,k,\alpha$ when those are known by context. 
In the next two lemmas, we show that finding  $S \in \mathcal{L}_{n,k,\alpha}(\Pi)$ (with appropriate parameters) is equivalent to finding~$S \in \cI_{n,k,\alpha}(\Pi)$ such that $w_n(S) = B_{n,k,\alpha}$. For the first direction, we have
\begin{lemma}
	\label{lemma:if}
	For every $n,k,\alpha \in \mathbb{N}_{>0}, \Pi \subseteq 2^{[n]}$, and $S \in \mathcal{L}_{n,k,\alpha}(\Pi)$, it holds that $w_{n}(S) = B_{n,k,\alpha}$. 
\end{lemma}
\begin{proof}
	Let $S \in \mathcal{L}(\Pi)$. By \eqref{eq:1}, $|S| = k$ and $\textnormal{\textsf{sum}}(S) = \alpha$. %In addition, let $D = \{x_i~|~i \in S\}$. 
	Therefore, \begin{equation*}
		\label{eq:ole}
		\begin{aligned}
			w_{n}(S) = \sum_{i \in S} \left( i+2 \cdot \textnormal{\textsf{sum}}([n])  \right) = \textnormal{\textsf{sum}}(S)+|S| \cdot 2 \cdot \textnormal{\textsf{sum}}([n]) = \alpha+2 k\cdot \textnormal{\textsf{sum}}([n]) = B_{n,k,\alpha}. 
		\end{aligned}
	\end{equation*}
	%	 The last equality holds because $\alpha = \textnormal{\textsf{sum}}(S) \leq \textnormal{\textsf{sum}}([n])$. 
	%	Since $S \in \mathcal{L}(\Pi)$ it holds that $S \in \cI$. Thus, by \eqref{eq:ole} $S$ is a solution for $U$ of profit exactly $B$. 
\end{proof}


For the second direction, we show the next result.
\begin{lemma}
	\label{lemma:onlyif}
	Let $n,k, \in \mathbb{N}_{>0}, \alpha \in \left[ n^2\right], \Pi \subseteq 2^{[n]}$, and $S \in \cI_{n,k,\alpha}(\Pi)$ such that $w_{n}(S) = B_{n,k,\alpha}$. Then, $S \in \mathcal{L}_{n,k,\alpha}(\Pi)$. 
\end{lemma}
\begin{proof}
	%Let $F$ be a solution for $U$ of profit at least $B$. Since $c(e) = p(e)$ for every element $e \in E$, then $p(F) = c(F) \leq B$, since $F$ satisfies the budget constraint; thus, since $p(F) \geq B$ we conclude that $p(F) = c(F) = B$. 
	Since $S \in \cI_{n,k,\alpha}(\Pi)$, by \eqref{eq:1}, it suffices to show that $|S| = k$ and $\textsf{sum}(S) = \alpha$. First, assume towards contradiction that $|S| \neq k$. By \eqref{eq:1}, this implies that  $|S|< k$. Thus,
	%We consider two cases. 
%	\begin{enumerate}
%		\item $|S|< k$. Then, 
\begin{equation}
			\label{eq:ol3}
			\begin{aligned}
				w_{n}(S) ={} & \sum_{i \in S} i+|S| \cdot 2 \cdot \textnormal{\textsf{sum}}([n]) \\
				\leq{} & \sum_{i \in S} i+(k-1) \cdot 2 \cdot \textnormal{\textsf{sum}}([n]) \\
				\leq{} &  \textnormal{\textsf{sum}}([n])+(k-1) \cdot 2 \cdot \textnormal{\textsf{sum}}([n]) \\
				<{} & 2 k\cdot \textnormal{\textsf{sum}}([n]) \\
				<{} & B_{n,k,\alpha}.
			\end{aligned}
		\end{equation} 
		
The last inequality holds since $\alpha \geq 1$. By \eqref{eq:ol3}, we reach a contradiction to $w_{n}(S) = B_{n,k,\alpha}$.		
		
%		\item $|S|>k$. Since $S \in \cI_{n,k,\alpha}(\Pi)$, by \eqref{eq:1} it holds that $|S| \leq k$ and we reach a contradiction. 
		
		\comment{
			\begin{equation}
				\label{eq:ol2}
				\begin{aligned}
					w_{n}(S) ={} & \sum_{i \in S} i+|S| \cdot 2 \cdot \textnormal{\textsf{sum}}([n]) \\
					\geq{} & \sum_{i \in S} i+(k+1) \cdot 2 \cdot \textnormal{\textsf{sum}}([n]) \\ 
					\geq{} & 2  \cdot \textnormal{\textsf{sum}}([n]) +2 k\cdot \textnormal{\textsf{sum}}([n]) \\
					>{} & B_{n,k,\alpha}.
				\end{aligned}
			\end{equation}
			The last inequality holds since: if $n = 1$ then $\alpha \leq 1 < 2 \cdot \textsf{sum}([n])$; otherwise, $n > 1$ and it follows that $\alpha \leq n^2 < n^2+n = 2 \cdot \textnormal{\textsf{sum}}([n])$. }%By \eqref{eq:ol2} we reach a contradiction since $F$ is a solution for $U$.
%	\end{enumerate}
We conclude that $|S| = k$. Therefore, 
	$$\textnormal{\textsf{sum}}(S) = \sum_{i \in S} i = w_{n}(S) - |S| \cdot 2 \cdot \textnormal{\textsf{sum}}([n]) = w_{n}(S) -  2 k \cdot \textnormal{\textsf{sum}}([n]) = B_{n,k,\alpha} - 2 k \cdot \textnormal{\textsf{sum}}([n]) = \alpha.$$ 
	In summary, we have $S \in \cI(\Pi)$, $|S| = k$ and $\textnormal{\textsf{sum}}(S) = \alpha$. By \eqref{eq:1}, it follows that $S \in \mathcal{L}(\Pi)$.
\end{proof}

The next result follows directly from \Cref{lemma:if,lemma:onlyif}.
\begin{cor}
	\label{cor:iff}
	For every $n,k \in \mathbb{N}_{>0}$, $\alpha \in \left[ n^2 \right]$, $ \Pi \subseteq 2^{[n]}$, and $S \in \mathcal{I}_{n,k,\alpha}(\Pi)$ it holds that $S \in \mathcal{L}_{n,k,\alpha}(\Pi)$ if and only if %there is $S \in \cI_{n,k,\alpha}(\Pi)$ such that 
	$w_{n}(S) = B_{n,k,\alpha}$. 
\end{cor}

We use \Cref{cor:iff} to derive lower bounds for several weighted matroid problems. For each of these problems, $\cP$, we use a reduction from $\Pi$-matroid. Given an instance $I$ of $\Pi$-matroid, we construct an instance $U$ of $\cP$ using the standard weight and standard budget.
The optimal solution for $U$ is {\em exactly} $B_{n,k,\alpha}$ (with appropriate parameters) iff $I$ has a solution. 
%\begin{proof}
%	%Clearly, if $\alpha \notin \left[ \textnormal{\textsf{sum}}([n])\right]$ then for all $S \subseteq [n]$ it holds that $\textsf{sum}(S) \neq \alpha$; thus, $\mathcal{L}_{n,k,\alpha}(\Pi) = \emptyset$. Therefore,  by \Cref{lemma:onlyif} there is no $F \in \cI_{n,k,\alpha}(\Pi)$ such that $w_{n}(F) = B_{n,k,\alpha}$ and the claim follows. Otherwise, $\alpha \notin \left[ \textnormal{\textsf{sum}}([n])\right]$. In this case, the proof follows by \Cref{lemma:if} and \Cref{lemma:onlyif}. 
%\end{proof}

%\begin{enumerate}
%	%\item There are $n$ elements $E = \{e_1, \ldots, e_n\}$.
%%	\item $\cI = \cI_{n,k,\alpha}(\Pi)$. 
%	\item For every $i \in [n]$ define $c(i) = p(i) = i+2 \cdot \textnormal{\textsf{sum}}([n])$. 
%	\item Define the budget as $B = \alpha+2 k \cdot \textnormal{\textsf{sum}}([n])$.
%\end{enumerate}


\subsection{Budgeted Matroid Independent Set}
\label{sec:abs_BM}

Recall that definition of the BM problem
\problem{ {\sc Budgeted Matroid Independent Set} (BM)}{
	{\bf Instance} &  
	$I = (E,\cI,c,p,B)$, where $(E,\cI)$ is a matroid, $c:E \rightarrow \mathbb{N}$ is  a cost function, $p:E \rightarrow \mathbb{N}$ is a  profit function and $B \in \mathbb{N}$ is the budget\\
	{\bf Solution}& $S\in \cI$ such that $c(S) \leq B$.\\
	{\bf Objective} & Find a solution $S$ of maximum profit $p(S)$. \\
	{\bf Model } &The collection of  independent sets $\cI$ is accessible only via a membership oracle; the encoding size of the instance $I$ is the encoding size of $(E,c,p,B)$ - the instance without $\cI$. 
	%We decided to repeat the model for all problems. --H 
	%\ariel{we need to re-discuss whether we want to repeat this definition over and over}
}
We first show how a reduction from the $\Pi$-matroid problem can be used to rule out the existence of a Fully PTAS for $\BM$. 

%\thmbm*

\begin{proof}
	
	
	%	\noindent{\bf Proof of \Cref{lem:hard}:}
	Assume towards a contradiction that there is a randomized Fully PTAS $\cA$ for $\BM$. We use $\cA$ to decide the $\Pi$-matroid problem.  Let $I = \left(n,k,\alpha,\mm\right)$ be a $\Pi$-matroid instance. Consider the following algorithm $\cB$ that decides $I$.
	\begin{enumerate}
		\item Construct the $\BM$ instance $%U_{n,k,\alpha}(\Pi) 
		U= ([n],\cI_{n,k,\alpha}(\Pi),w_{n},w_{n},B_{n,k,\alpha})$ along with the membership oracle of $M_{n,k,\alpha}(\Pi)$; note that the elements of $U$ are $[n] = \{1,\ldots, n \}$, the independent sets are the independent sets of the $\Pi$-matroid $\mm$ (i.e., $\cI_{n,k,\alpha}(\Pi)$), the standard weight function $w_n$ is both the profit and the cost function, and the budget is simply the standard budget with parameters $n,k,\alpha$. %Let $U = U_{n,k,\alpha}(\Pi)$ for simplicity.
		\item Define $\eps = \frac{1}{6 \cdot n^3}$ and compute $\cA$ on $U$ and $\eps$; let $S$ be the returned solution.  
		\item Return that $I$ has a solution if and only if $w_n(S) = B_{n,k,\alpha}$. 
	\end{enumerate}
	
	Note that $U$ can be constructed from $I$ in time $\io$ and that $\frac{1}{\eps} = |I|^{O(1)}$; thus, the running time of $\cB$ on $I$ is $|I|^{O(1)}$, since $\cA$ is a (randomized) FPTAS. We now show correctness. %For the correctness, %we show that $\cB$ returns that $I$ has a solution if and only if $I$ indeed has a solution. 
	\begin{itemize}
		\item If $\cB$ returns that $I$ has a solution. Then, $w_n(S) = B_{n,k,\alpha}$; since $S$ is a solution for $U$ it also holds that $S \in \cI_{n,k,\alpha}(\Pi)$. Thus, by \Cref{cor:iff} it holds that $I$ has a solution.  
		\item If $I$ has a solution. Then, there is $T \in \mathcal{L}_{n,k,\alpha}(\Pi)$; thus, by  \Cref{cor:iff} it follows that $w_n(T) = B_{n,k,\alpha}$ implying that $T$ is a solution for $U$.  Therefore, with probability at least $\frac{1}{2}$: 
		 \begin{equation}
			\label{eq:sOPT}
			B_{n,k,\alpha} \geq \OPT(U) \geq w_n(S) \geq (1-\eps) \cdot w_n(T) \geq  w_n(T)-\frac{1}{6 \cdot n^3} \cdot 3 \cdot n^3 = B_{n,k,\alpha}-\frac{1}{2}.
		\end{equation} 	The first inequality holds since %$w_n(T) = B_{n,k,\alpha}$. In addition, $\OPT(U) \leq B_{n,k,\alpha}$ because 
	the profit and the cost of each element are the same, and the cost of a solution cannot exceed the budget. The third inequality follows (with probability at least $\frac{1}{2}$) because $S,T$ are solutions for $U$, and $S$ is computed by a randomized FPTAS with the parameter $\eps$. The fourth inequality follows from \Cref{lem:OPTu} and by the selection of $\eps$. Observe that the value of $S$ is an integer by \eqref{eq:w}; thus, by \eqref{eq:sOPT}, %w.p. at least $\frac{1}{2}$ 
it holds that $w_n(S) = B_{n,k,\alpha}$ w.p. at least $\frac{1}{2}$. Then, $\cB$ returns that $I$ has a solution with probability at least $\frac{1}{2}$.  
	\end{itemize}
	%
	%by \Cref{cor:iff} a solution $T$ for $U$  is a solution for $I$ if and only if %therefore, by  it holds that %algorithm $\cA$ returns a solution for $U$ with profit $B_{n,k,\alpha}$

	
	
	%Therefore, there is a solution for $I$ if and only if $\OPT(U) = B_{n,k,\alpha}$; thus, to conclude we show that $w_n(S) = \OPT(U)$ w.p. at least $\frac{1}{2}$. %if and only if $S \in \mathcal{L}_{n,k,\alpha}(\Pi)$ (i.e., $I$ has a solution).
	%	  
	%  
	%	  we show that w.p. at least $\frac{1}{2}$ it holds that: $w_n(S) = B_{n,k,\alpha}$ if and only if I has a solution. 
%	By the approximation guarantee of $\cA$ with our parameter $\eps$, w.p. at least $\frac{1}{2}$ the solution $S$ satisfies \begin{equation}
%		\label{eq:sOPT}
%		\OPT(U)-w_n(S) \leq \eps \cdot \OPT(I) \leq  \frac{1}{6 \cdot n^3} \cdot 3 \cdot n^3 = \frac{1}{2}.
%	\end{equation} The first inequality follows because $S$ is a solution computed by an FPTAS with the parameter $\eps$. The second inequality holds since $\OPT(U) \leq 3 \cdot n^3$ from \Cref{lem:OPTu} and by the selection of $\eps$. Observe that the value of $S$ is an integer by \eqref{eq:w}; thus, by \eqref{eq:sOPT}, w.p. at least $\frac{1}{2}$ it holds that $w_n(S) = \OPT(U)$.  %, by the definition of a randomized FPTAS and since $\eps^{-1} = 2 \cdot |I|^c = \io$.  
	%  
	%If $\alpha > \textnormal{\textsf{sum}}([n])$ by \Cref{def:Matroid} and \Cref{def:problem} there is no solution for $I$ and deciding $I$ is trivial. Thus, for the following, assume w.l.o.g. that $\alpha \leq \textnormal{\textsf{sum}}([n])$. 
	%We define the BM instance $U = (E, \cI, c, p, B)$ such that the following holds. 
	%
	%
	%	By \Cref{lem:OPTu}, for all $S \subseteq [n]$ it holds that $w_{n}(S) \in \mathbb{N}$ and that $\OPT \left( U\right) = n^{O(1)} = \left| U \right|^{O(1)}$; thus, $U$ is a scaled instance of BM. Therefore, if BM has a randomized FPTAS, by \Cref{lem:FPTAStoEXACT} there is a randomized exact algorithm $\cA$ that finds in time $|U|^{O(1)} = \io$ an optimal solution for $U$. 
	%	Clearly, $\OPT(U) \leq B_{n,k,\alpha}$ because the profit and the cost of each element are the same, and the cost of a solution cannot exceed the budget. Since $S$ is a solution for $U$ it follows that $S \in \cI_{n,k,\alpha}(\Pi)$; therefore, by \Cref{cor:iff} it holds that %algorithm $\cA$ returns a solution for $U$ with profit $B_{n,k,\alpha}$
	%	Thus, w.p. at least $\frac{1}{2}$, it holds that $w_n(S) = B_{n,k,\alpha}$ if and only if $S \in \mathcal{L}_{n,k,\alpha}(\Pi)$ (i.e., $I$ has a solution). 
	Hence %by computing $\cA$ on $U$ and getting a solution $S$, 
	the algorithm  correctly decides $I$ with probability at least $\frac{1}{2}$. %if and only if $w_n(S) = \OPT(U)$, which holds w.p. at least $\frac{1}{2}$. 
	It follows that $\cB$ is %by returning that $I$ has a solution if and only if $w_n(S) = B_{n,k,\alpha}$. Thus, we 
	a randomized algorithm which decides the $\Pi$-matroid problem, contradicting \Cref{cor:deter}. 
\end{proof}



\subsection{Exact Weight Matroid Independent Set}
\label{sec:abs_EMI}




%	We say that a problem is {\em Strongly} NP-Hard if it remains NP-Hard if all numbers in the instance of the problem are encoded in unary. 
The proof of \Cref{thm:emi} is similar to the proof of \Cref{thm:bm}.
Recall that $S$ is a {\em solution} for an EMI instance $(E,\cI,w,k,T)$ if $S \in \cI$, $|S| = k$, and $w(S) = T$. 
% \begin{lemma}
%	\label{lem:EMI}
%	For any function $f:\mathbb{N} \rightarrow \mathbb{N}$ there is no algorithm that decides \textnormal{EMI} in time $f(k) \cdot (T \cdot n)^{O(1)}$, where $n$ is the number of elements in the given matroid, $k$ is the rank of the matroid, and $T$ is the target value.  
%%	Unless \textnormal{P=NP}, there is no polynomial time algorithm that decides \textnormal{EMI}. %is \textnormal{Strongly NP-Hard}.  
%\end{lemma}

%\thmemi*

\begin{proof}
	Assume towards contradiction that there is a function $f:\mathbb{N} \rightarrow \mathbb{N}$,  and an algorithm $\cA$ for  \textnormal{EMI} which, given an instance $I=(E,\cI,w,k,T)$, runs in time  $f(k) \cdot \left(T \cdot |I|\right)^{O(1)}$.  
	%We use $\cA$ to attain a randomized algorithm for  the $\Pi$-matroid problem. 
	%By \Cref{lem:randomized}, there is a constant $c$ such that there is no randomized algorithm that decides the $\Pi$-matroid problem in time $c \cdot \frac{{n' \choose k}}{{n'}^3}$. %randomized FPTAS $\cA$ for EMI. 
 % Let $I = \left(n,k,\alpha,\mm\right)$ be a $\Pi$-matroid instance such that $7 \cdot n^3+f(k) \cdot (2 \cdot n^4)^{q} < c \cdot \frac{{n \choose k}}{{n}^3}$. Clearly, since $c,q$ are constants, the above inequality is realized by some $n,k$, where $n$ is sufficiently large w.r.t. $k$.
   Consider the following algorithm for the $\Pi$-matroid problem whose input is instance $I=(n,k,\alpha,M_{n,k,\alpha}(\Pi))$, to which we refer as  $\cB$, that decides $I$.
	\begin{enumerate}
		\item Construct the EMI instance $%U_{n,k,\alpha}(\Pi) 
		U= ([n],\cI_{n,k,\alpha}(\Pi),w_{n},k,B_{n,k,\alpha})$ along with the membership oracle of $M_{n,k,\alpha}(\Pi)$; note that the elements of $U$ are $[n] = \{1,\ldots, n \}$, the independent sets are the independent sets of the $\Pi$-matroid $\mm$ (i.e., $\cI_{n,k,\alpha}(\Pi)$), the standard weight function $w_n$ is the weight function, the target cardinality is $k$, and the target value is the standard budget with the parameters $n,k,\alpha$. %We say that $T \subseteq [n]$ is a solution for $U$ if $T \in \cI_{n,k,\alpha}$ and $w_n(T) = B_{n,k,\alpha}$. %Let $U = U_{n,k,\alpha}(\Pi)$ for simplicity.
		\item %Define $\eps = \frac{1}{6 \cdot n^3}$ and 
		Execute $\cA$ on the instance  $U$.  %and let $S$ be the returned solution.  
		 Return that $I$ has a solution if and only if $\cA$ returned that $U$ has a solution.  %$w_n(S) = B_{n,k,\alpha}$. 
	\end{enumerate}
	
	We first assert that $\cB$ is a randomized algorithm for the $\Pi$-matroid problem.
		\begin{itemize}
		\item If $\cB$ returns that $I$ has a solution, then the EMI instance $U$ has a solution $S$. 
		That is,~$w_n(S) = B_{n,k,\alpha}$ and $S \in \cI_{n,k,\alpha}(\Pi)$. Thus, by \Cref{cor:iff} it holds that $S\in \mathcal{L}_{n,k,\alpha}(\Pi)$. That is, $S$ is a solution for $I$, and the response of the algorithm is correct. 
		\item If $I$ has a solution then there is $T \in \mathcal{L}_{n,k,\alpha}(\Pi)$; thus, by  \Cref{cor:iff}, it holds that  $w_n(T) = B_{n,k,\alpha}$ and $|T|=k$ by \eqref{eq:1}, implying that $T$ is a solution for $U$. Therefore, with probability  at least $\frac{1}{2}$ it holds that $\cA$ returns that $U$ has a solution; thus, with probability  at least $\frac{1}{2}$ it holds that $\cB$ returns that $I$ has a solution. 
		%		\begin{equation}
			%			\label{eq:sOPT}
			%			B_{n,k,\alpha} \geq \OPT(U) \geq w_n(S) \geq (1-\eps) \cdot w_n(T) \geq  w_n(T)-\frac{1}{6 \cdot n^3} \cdot 3 \cdot n^3 = B_{n,k,\alpha}-\frac{1}{2}.
			%		\end{equation} 	The first inequality holds since %$w_n(T) = B_{n,k,\alpha}$. In addition, $\OPT(U) \leq B_{n,k,\alpha}$ because 
		%		the profit and the cost of each element are the same, and the cost of a solution cannot exceed the budget. The third inequality follows (w.p. at least $\frac{1}{2}$) because $S,T$ are solutions for $U$, and $S$ is computed by a randomized FPTAS with the parameter $\eps$. The fourth inequality follows from \Cref{lem:OPTu} and by the selection of $\eps$. Observe that the value of $S$ is an integer by \eqref{eq:w}; thus, by \eqref{eq:sOPT}, %w.p. at least $\frac{1}{2}$ 
		%		it holds that $w_n(S) = B_{n,k,\alpha}$ w.p. at least $\frac{1}{2}$. Then, $\cB$ returns that $I$ has a solution w.p. at least $\frac{1}{2}$.  
	\end{itemize}
	Hence, 
	the algorithm correctly decides  $I$ with probability at least $\frac{1}{2}$. %if and only if $w_n(S) = \OPT(U)$, which holds w.p. at least $\frac{1}{2}$. 
	It follows that $\cB$ is %by returning that $I$ has a solution if and only if $w_n(S) = B_{n,k,\alpha}$. Thus, we 
	a randomized algorithm which decides the $\Pi$-matroid problem.
	
	
	
	
	Note that $B_{n,k,\alpha} = 2k \cdot \textsf{sum}(n) \leq 2k \cdot n^2 \leq 2 \cdot n^3$ since $k \leq n$. 
	Thus, $U$ can be constructed in time $|I|^{O(1)}$. 
	Let $q(n,k)$ be the maximum running time of $\cB$ on an instance of the $\Pi$-matroid problem in which the number of elements is $n$ and  the parameter is $k$. Since $\cA$ runs in time $f(k)\cdot (T\cdot |I|)^{O(1)}$ it follows that there is $q\in \mathbb{N}$ such that 
	 $q(n,k)\leq f(k)\cdot (n\cdot 2\cdot n^3)^{O(1)} \leq f(k)\cdot n^{q}$.
	  Furthermore, by  \Cref{lem:randomized} it also holds that $q(n,k) =\Omega\left( \frac{ {n \choose k }}{n^3}\right) $. 
	  Thus, there are $c>0$ and $M>0$ such that for every $n,k>M$ it holds that $q(n,k) \geq c\cdot \frac{ {n \choose k }}{n^3}$. 
	  Select $\tilde{k}=\max\left\{ M+1, q+10\right\}$, then for every $n>M$ such that $k$ divides $n$ we have
	 $$
	 f(\tilde{k})\cdot n^{q} \geq q(n,\tilde{k})\geq  c\cdot \frac{ {n \choose \tilde{k}}}{n^3} \geq c\cdot \frac{\left(\frac{n}{\tilde{k}}\right)^{\tilde{k}}}{n^3 } ,
	 $$
	 where the last inequality uses the inequality ${ n\choose k } \geq \left(\frac{n}{k} \right)^k$ when $k$ divides $n$.
	 Thus,
	 $$
	 f(\tilde{k})\cdot \tilde{k}^{\tilde{k}} \geq   c\cdot n^{\tilde{k}-3-q}.
	 $$
	 The left hand term in the above equation is a constant  in $n$, while the right hand side goes to infinity as $n$ goes to infinity. Thus, we have a contradiction. Therefore, there is no randomized algorithm for EMI which runs in time $f(k)\cdot (T\cdot |I|)^{O(1)}$.
	
	
	%$7 \cdot n^3+f(k) \cdot (B_{n,k,\alpha} \cdot n)^{q} \leq 7 \cdot n^3+f(k) \cdot (2 \cdot n^4)^{q}$. We now show correctness.  %For the correctness, %we show that $\cB$ returns that $I$ has a solution if and only if $I$ indeed has a solution. 
 
\end{proof}


\comment{
	
	\begin{proof}
		
		
		
		We use a reduction from the $\Pi$-matroid problem.	Let $I = \left(n,k,\alpha,\mm\right)$ be a $\Pi$-matroid instance. We define the following BM instance $U = ([n], \cI, w, T)$ such that the following holds. \begin{enumerate}
			%\item There are $n$ elements $E = \{e_1, \ldots, e_n\}$.
			\item $\cI = \cI_{n,k,\alpha}(\Pi)$. 
			\item Define a weight function $w:[n] \rightarrow \mathbb{R}_{\geq 0}$, where for all $vi \in [n]$ it holds that $w(i) = 1+\frac{i}{2 \cdot n^2}$.  
			\item $T = k+\frac{\alpha}{2 \cdot n^2}$. 
			%			\begin{equation}
				%				\label{eq:w}
				%				w(v) = 1+\frac{f(v)}{2 \cdot n^2}. 
				%			\end{equation}
		\end{enumerate}  
		
		%	Let $X = (E,v,\sigma,D,T)$ be an SSG instance and let $n = |E|$. For every $k \in \{0,1\ldots,n\}$ 
		
		%		Define the matroid $M_{\alpha} = (V,\cI_{\alpha})$, $\cI_{\alpha} = \cI^1_{\alpha} \cup \cI^2_{\alpha} \cup \cI^3_{\alpha}$, such that \begin{equation}
			%			\label{eq:1}
			%			\begin{aligned}
				%				\cI^1_{\alpha} = {} & \{S \subseteq V~\big|~ |S| < k\} \\
				%				\cI^2_{\alpha} = {} & \{S \subseteq V~\big|~ |S| = k, f(S) \neq \alpha\} \\
				%				\cI^3_{\alpha} = {} & \{S \subseteq V~\big|~ |S| = k, f(S) = \alpha, S \in C(G)\}. \\
				%			\end{aligned}
			%		\end{equation}
		
		
		%	For every $\alpha \in [n^2]$ let $T_{\alpha} = k+\frac{\alpha}{2 \cdot n^2}$ and define the EMI instance $I_{\alpha} = (V,\cI_{\alpha},w,T_{\alpha})$.
		
		
		% such that the following holds. \begin{equation}
			%	\label{eq:sB2}
			%	\begin{aligned}
				%	c(e) ={} & \beta \cdot \sigma(e) +v(e)  ~~\forall e \in E \\
				%	p(e) ={} & \alpha+\sigma(e)~~\forall e \in E \\
				%	B ={} & \beta \cdot D+T.
				%	\end{aligned}
			%\end{equation}
			%
			
			\begin{claim}
				\label{claim:r<4}
				For every $\alpha \in [n^2]$ and $Q \subseteq V$ it holds that $Q$ is a solution for $I_{\alpha}$ if and only if $Q \in C(G)$, $|Q| = k$, and $f(Q) = \alpha$. 
			\end{claim}
			
			\begin{proof}
				Let $Q$ be a solution for $I_{\alpha}$. Then, \begin{equation}
					\label{eq:2}
					|Q| \leq w(Q) = T_{\alpha} = k+\frac{\alpha}{2 \cdot n^2} \leq k+\frac{n^2}{2 \cdot n^2} < k+1.  
				\end{equation}
				and 
				\begin{equation}
					\label{eq:3}
					|Q| = w(Q)- \frac{f(Q)}{2 \cdot n^2} = T_{\alpha} - \frac{f(Q)}{2 \cdot n^2}= k+\frac{\alpha}{2 \cdot n^2}- \frac{f(Q)}{2 \cdot n^2} \geq k- \frac{f(Q)}{2 \cdot n^2} \geq k - \frac{n^2}{2 \cdot n^2} > k-1.  
				\end{equation} The second inequality holds since $f(V) \leq n \cdot \max_{v \in V} f(V) = n^2$. By \eqref{eq:2} and \eqref{eq:3} we conclude that $|Q| = k$. Moreover, %it holds that $f(Q) = \frac{\alpha}{}$
				\begin{equation}
					\label{eq:a}
					f(Q) = (w(Q) - k) \cdot 2 n^2  = (T_{\alpha}-k) \cdot 2 n^2 = \alpha.  
				\end{equation} The first equality follows from \eqref{eq:w}. We get that $f(Q) = \alpha$ by \eqref{eq:a}. Therefore, since $|Q| = k$, $f(Q) = \alpha$, and $Q \in \cI_{\alpha}$ ($Q$ is a solution for $I_{\alpha}$), it follows that  $Q \in \cI^3_{\alpha}$; thus, $Q \in C(G)$ by \eqref{eq:1}. 
				
				For the other direction, let $Q \subseteq V$  such that $Q \in C(G)$, $|Q| = k$, and $f(Q) = \alpha$. Then, by \eqref{eq:1} it holds that $Q \in \cI^3_{\alpha}$ and therefore $Q \in \cI_{\alpha}$. Moreover, by \eqref{eq:w} it holds that $w(Q) = k+\frac{f(Q)}{2n^2} = k+\frac{\alpha}{2 \cdot n^2} = T_{\alpha}$.  We conclude that $Q$ is a solution for $I_{\alpha}$. 
			\end{proof}
			
			
			To prove the lemma, consider the following cases. \begin{enumerate}
				\item There is a clique $Q \in C(G)$ of size $|Q| = k$. Let $\alpha = f(Q)$. Then, by Claim~\ref{claim:r<4} it holds that $Q$ is a solution for $I_{\alpha}$. 
				
				\item  There is no clique of size $k$ in $G$. Then, assume towards a contradiction that there is $\alpha \in [n^2]$ such that there is a solution $Q$ for $I_{\alpha}$. By Claim~\ref{claim:r<4} it holds that $Q$ is in particular a clique of size $k$ in $G$ in contradiction. Hence, for all $\alpha \in [n^2]$ there is no solution for $I_{\alpha}$ in this case. 
			\end{enumerate} Thus, by the above, to decide if there is a clique of size $k$ in $G$ we can decide if there is a solution for at least one $I_{\alpha}$ for all $\alpha \in [n^2]$. The reduction takes polynomial time and the weights $w$ are polynomial (in the encoding size of the $k$-clique instance). The proof follows since the $k$-clique problem is known to be strongly NP-Hard.  		\end{proof}
		
		
	}
	
	
	\subsection{Constrained Minimum Basis of a Matroid}
	\label{sec:abs_CMB}
%		\begin{definition}[Constrained Minimum Basis of a Matroid (BM)] 
%		\label{def:CMB}We are given a tuple $U = (E,\cI,\ell,w,L)$, where $(E,\cI)$ is a matroid along with a membership oracle, $c:E \rightarrow \mathbb{N}$ is a cost function, $w:E \rightarrow \mathbb{N}$ is a weight function and $L \in \mathbb{N}$ is a bound. A solution of $U$ is a basis $S\in \cI, |S| = \max_{S' \in \cI} |S'|$ such that $\ell(S) \leq L$. The goal is to find a solution $S$ of $U$ with minimum weight $w(S)$. 
%	\end{definition} 
	
	
	Unlike the reductions used from  BM and EMI,  the reduction from $\Pi$-matroid to CMB uses a  different {\em cost} function. Specifically, for every $n \in \mathbb{N}_{>0}$ define the standard cost function of $n$ as $c_n:[n] \rightarrow \mathbb{N}$ such that for all $i \in [n]$ it holds that \begin{equation}
		\label{leq:w}
		c_n(i) = w_n(i)-2 \cdot i = 2 \cdot \textsf{sum}([n])-i.  
	\end{equation} Moreover, define the standard bound of  $n,k,\alpha \in \mathbb{N}_{>0}$  as \begin{equation}
		\label{leq:B}
		L_{n,k,\alpha} = B_{n,k,\alpha} -2 \cdot \alpha =2 k\cdot \textnormal{\textsf{sum}}([n])- \alpha.   
	\end{equation}
%The above defined cost function is useful for the proof of hardness of CMB. We first use the following auxiliary results.  
	
	
%		\begin{lemma}
%		\label{lemma:wl}
%		For every $n,k,\alpha \in \mathbb{N}_{>0}$, and $S \subseteq [n], |S| = k$ such that $w_n(S) = B_{n,k,\alpha}$, it holds that $c_n(S) \leq L_{n,k,\alpha}$. 
%	\end{lemma}
%	
%	\begin{proof}
%	 By \eqref{eq:w}, \eqref{eq:B}, and that $|S| = k$ we have: \begin{equation}
%		\label{eq:x}
%		2 \cdot \textsf{sum}([n]) \cdot k+\alpha = B_{n,k,\alpha} = w_n(S) = 2 \cdot \textsf{sum}([n]) \cdot k+\textsf{sum}(S).
%	\end{equation} By \eqref{eq:x} we conclude that $\textsf{sum}(S) = \alpha$. Therefore, 
%	
%		$$c_n(S) = 2 \cdot \textsf{sum}([n]) \cdot k-\textsf{sum}(S) = 2 \cdot \textsf{sum}([n]) \cdot k-2\cdot \alpha = L_{n,k,\alpha}.$$
%	\end{proof}

We use the next auxiliary lemma in the proof of \Cref{thm:CMB}. 
	\begin{lemma}
	\label{lemma:wl}
	For every $n,k,\alpha \in \mathbb{N}_{>0}$, and $S \subseteq [n]$ such that  $|S| = k$ it holds that $w_n(S) \geq B_{n,k,\alpha}$ if and only if $c_n(S) \leq L_{n,k,\alpha}$. 
\end{lemma}

\begin{proof}
	If $w_n(S) \geq B_{n,k,\alpha}$, then by \eqref{eq:w}, \eqref{eq:B}, and since $|S| = k$ we have: \begin{equation}
		\label{eq:x}
		2 \cdot \textsf{sum}([n]) \cdot k+\alpha = B_{n,k,\alpha} \leq w_n(S) = 2 \cdot \textsf{sum}([n]) \cdot k+\textsf{sum}(S).
	\end{equation} By \eqref{eq:x} we conclude that $\textsf{sum}(S) \geq \alpha$. Therefore, 
	$$c_n(S) = 2 \cdot \textsf{sum}([n]) \cdot k-\textsf{sum}(S) \leq 2 \cdot \textsf{sum}([n]) \cdot k- \alpha = L_{n,k,\alpha}.$$
	For the second direction, if $c_n(S) \leq L_{n,k,\alpha}$, then by \eqref{leq:w}, \eqref{leq:B}, and since $|S| = k$ it holds that: \begin{equation}
		\label{leq:x}
		2 \cdot \textsf{sum}([n]) \cdot k-\alpha = L_{n,k,\alpha} \geq c_n(S) = 2 \cdot \textsf{sum}([n]) \cdot k-\textsf{sum}(S).
	\end{equation} By \eqref{leq:x} we conclude that $\textsf{sum}(S) \geq \alpha$. Then, 
	$$w_n(S) = 2 \cdot \textsf{sum}([n]) \cdot k+\textsf{sum}(S) \geq 2 \cdot \textsf{sum}([n]) \cdot k+ \alpha = B_{n,k,\alpha}.$$
\end{proof}

%	\begin{lemma}
%		\label{llemma:if}
%		For every $n,k,\alpha \in \mathbb{N}_{>0}, \Pi \subseteq 2^{[n]}$, and $S \in \mathcal{L}_{n,k,\alpha}(\Pi)$, it holds that $c_n(S) = L_{n,k,\alpha}$. 
%	\end{lemma}
%	\begin{proof}
%		Let $S \in \mathcal{L}(\Pi)$. That is, by \eqref{eq:1} it holds that $|S| = k$ and $\textnormal{\textsf{sum}}(S) = \alpha$. %In addition, let $D = \{x_i~|~i \in S\}$. 
%		Then, \begin{equation*}
%			\label{leq:ole}
%			\begin{aligned}
%				c_n(S) = \sum_{i \in S} \left( 2 \cdot \textnormal{\textsf{sum}}([n])-i  \right) = |S| \cdot 2 \cdot \textnormal{\textsf{sum}}([n])-\textnormal{\textsf{sum}}(S) = 2 k\cdot \textnormal{\textsf{sum}}([n])-\alpha = L_{n,k,\alpha}.  
%			\end{aligned}
%		\end{equation*}
%		%	 The last equality holds because $\alpha = \textnormal{\textsf{sum}}(S) \leq \textnormal{\textsf{sum}}([n])$. 
%		%	Since $S \in \mathcal{L}(\Pi)$ it holds that $S \in \cI$. Thus, by \eqref{eq:ole} $S$ is a solution for $U$ of profit exactly $B$. 
%	\end{proof}
	
	
%	Now, for the second direction, given $S \in \cI_{n,k,\alpha}(\Pi)$ such that $c_n(S) = B_{n,k,\alpha}$ we show that $S \in \mathcal{L}_{n,k,\alpha}(\Pi)$.  
%	
%	\begin{lemma}
%		\label{llemma:onlyif}
%		Let $n,k, \in \mathbb{N}_{>0}, \alpha \in \left[ n^2\right], \Pi \subseteq 2^{[n]}$, and $S \in \cI_{n,k,\alpha}(\Pi)$ such that $w_{n}(S) = B_{n,k,\alpha}$. Then, $S \in \mathcal{L}_{n,k,\alpha}(\Pi)$. 
%	\end{lemma}
%	\begin{proof}
%		%Let $F$ be a solution for $U$ of profit at least $B$. Since $c(e) = p(e)$ for every element $e \in E$, then $p(F) = c(F) \leq B$, since $F$ satisfies the budget constraint; thus, since $p(F) \geq B$ we conclude that $p(F) = c(F) = B$. 
%		%For the following, 
%		%	We show that $F$ is also a solution for $I$. 
%		Since $S \in \cI_{n,k,\alpha}(\Pi)$, by \eqref{eq:1}, it is sufficient to show that $|S| = k$ and $\textsf{sum}(S) = \alpha$. First, assume towards a contradiction that $|S| \neq k$. We consider two cases. \begin{enumerate}
%			\item $|S|< k$. Then, \begin{equation}
%				\label{eq:ol3}
%				\begin{aligned}
%					w_{n}(S) ={} & \sum_{i \in S} i+|S| \cdot 2 \cdot \textnormal{\textsf{sum}}([n]) \\
%					\leq{} & \sum_{i \in S} i+(k-1) \cdot 2 \cdot \textnormal{\textsf{sum}}([n]) \\
%					\leq{} &  \textnormal{\textsf{sum}}([n])+(k-1) \cdot 2 \cdot \textnormal{\textsf{sum}}([n]) \\
%					<{} & 2 k\cdot \textnormal{\textsf{sum}}([n]) \\
%					<{} & B_{n,k,\alpha}.
%				\end{aligned}
%			\end{equation} 
%			
%			The last inequality holds since $\alpha \geq 1$. By \eqref{eq:ol3} we reach a contradiction that $w_{n}(S) = B_{n,k,\alpha}$.%	By the above we reach contradiction that $p(F) = B$. 
%			\item $|S|>k$. Since $S \in \cI_{n,k,\alpha}(\Pi)$, by \eqref{eq:1} it holds that $|S| \leq k$ and we reach a contradiction. 
%			
%			\comment{
%				\begin{equation}
%					\label{eq:ol2}
%					\begin{aligned}
%						w_{n}(S) ={} & \sum_{i \in S} i+|S| \cdot 2 \cdot \textnormal{\textsf{sum}}([n]) \\
%						\geq{} & \sum_{i \in S} i+(k+1) \cdot 2 \cdot \textnormal{\textsf{sum}}([n]) \\ 
%						\geq{} & 2  \cdot \textnormal{\textsf{sum}}([n]) +2 k\cdot \textnormal{\textsf{sum}}([n]) \\
%						>{} & B_{n,k,\alpha}.
%					\end{aligned}
%				\end{equation}
%				The last inequality holds since: if $n = 1$ then $\alpha \leq 1 < 2 \cdot \textsf{sum}([n])$; otherwise, $n > 1$ and it follows that $\alpha \leq n^2 < n^2+n = 2 \cdot \textnormal{\textsf{sum}}([n])$. }%By \eqref{eq:ol2} we reach a contradiction since $F$ is a solution for $U$.
%		\end{enumerate}By the above we conclude that $|S| = k$. Therefore, 
%		$$\textnormal{\textsf{sum}}(S) = \sum_{i \in S} i = w_{n}(S) - |S| \cdot 2 \cdot \textnormal{\textsf{sum}}([n]) = w_{n}(S) -  2 k \cdot \textnormal{\textsf{sum}}([n]) = B_{n,k,\alpha} - 2 k \cdot \textnormal{\textsf{sum}}([n]) = \alpha.$$ In conclusion, we have $S \in \cI(\Pi)$, $|S| = k$ and $\textnormal{\textsf{sum}}(S) = \alpha$. Thus, by \eqref{eq:1} it holds that $S \in \mathcal{L}(\Pi)$.
%	\end{proof}

	
%	\begin{lemma}
%	\label{lem:L}
%	For every $n,k,\alpha \in \mathbb{N}_{>0}$, and $S \subseteq [n]$ it holds that $w_n(S)  \leq B_{n,k,\alpha}$ if and only if $c_n(S) \geq L_{n,k,\alpha}$. 
%\end{lemma}

%\begin{proof}
%	Assume that $w_n(S)  \leq B_{n,k,\alpha}$; then, $$c_n(S) = w_n(S) -2 \cdot \textsf{sum}(S) w_n(S) \geq L_{n,k,\alpha}$$
%\end{proof}
%\thmcmb*
	
	\begin{proof}
		Assume towards a contradiction that there is a randomized Fully PTAS $\cA$ for CMB. We use $\cA$ to decide the $\Pi$-matroid problem.  Let $I = \left(n,k,\alpha,\mm\right)$ be a $\Pi$-matroid instance. Consider the following algorithm $\cB$ that decides $I$.
		\begin{enumerate}
			\item Construct the CMB instance $%U_{n,k,\alpha}(\Pi) 
			U= ([n],\cI_{n,k,\alpha}(\Pi),c_n,w_{n},L_{n,k,\alpha})$ along with the membership oracle of $M_{n,k,\alpha}(\Pi)$; note that the elements of $U$ are $[n] = \{1,\ldots, n \}$, the independent sets are the independent sets of the $\Pi$-matroid $\mm$ (i.e., $\cI_{n,k,\alpha}(\Pi)$), the cost is the standard cost function, the weight is the standard weight function $w_n$, and the bound is the standard bound with parameters $n,k,\alpha$. %Let $U = U_{n,k,\alpha}(\Pi)$ for simplicity.
			\item Define $\eps = \frac{1}{6 \cdot n^3}$ and compute $\cA$ on $U$ and $\eps$; let $S$ be the returned solution.  
			\item Return that $I$ has a solution if and only if $w_n(S) = B_{n,k,\alpha}$. 
		\end{enumerate}
		Note that $U$ can be constructed from $I$ in time $\io$ and that $\frac{1}{\eps} = |I|^{O(1)}$; thus, the running time of $\cB$ on $I$ is $|I|^{O(1)}$, since $\cA$ is a (randomized) Fully PTAS. We now show correctness. %For the correctness, %we show that $\cB$ returns that $I$ has a solution if and only if $I$ indeed has a solution. 
		\begin{itemize}
			\item If $\cB$ returns that $I$ has a solution then  $w_n(S) = B_{n,k,\alpha}$; since $S$ is a solution for $U$ it also holds that $S \in \cI_{n,k,\alpha}(\Pi)$. Thus, by \Cref{cor:iff} it holds that $S\in \mathcal{L}_{n,k,\alpha}(\Pi)$, and thus  $I$ has a solution.  
			\item If $I$ has a solution then there is $T \in \mathcal{L}_{n,k,\alpha}(\Pi)$; thus, by  \Cref{cor:iff} it follows that $w_n(T) = B_{n,k,\alpha}$, which implies by \Cref{lemma:wl} that $c_n(T) \leq L_{n,k,\alpha}$. Then, since $T$ is a basis of $\mm$ ($T \in \mathcal{L}_{n,k,\alpha}(\Pi)$), it follows that $T$ is a solution for $U$.  Hence, with probability at least $\frac{1}{2}$:  
			\begin{equation}
				\label{eq:sOPTcmb}
				B_{n,k,\alpha} \leq \OPT(U) \leq w_n(S) \leq (1+\eps) \cdot w_n(T) \leq  w_n(T)+\frac{1}{6 \cdot n^3} \cdot 3 \cdot n^3 = B_{n,k,\alpha}+\frac{1}{2}.
			\end{equation} 	The first inequality follows from \Cref{lemma:wl} (note that $U$ has a solution since $T$ is a solution). %for every $S' \subseteq [n], |S'|<k$ it holds that $s$ every solution $S'$ for $U$ is a basis of $\mm$; thus, it follows that $w$ %$w_n(T) = B_{n,k,\alpha}$. In addition, $\OPT(U) \leq B_{n,k,\alpha}$ because 
		%	the profit and the cost of each element are the same, and the cost of a solution cannot exceed the budget. 
		The third inequality follows (with probability at least $\frac{1}{2}$) because $S,T$ are solutions for $U$, and $S$ is computed by a randomized Fully PTAS with the parameter $\eps$. The fourth inequality follows from \Cref{lem:OPTu} and by the selection of $\eps$. Observe that the value of $S$ is an integer by \eqref{eq:w}; thus, by \eqref{eq:sOPTcmb}, %w.p. at least $\frac{1}{2}$ 
			it holds that $w_n(S) = B_{n,k,\alpha}$ with probability  at least $\frac{1}{2}$. Thus, $\cB$ returns that $I$ has a solution w.p. at least $\frac{1}{2}$.  
		\end{itemize}
		Hence, %by computing $\cA$ on $U$ and getting a solution $S$, 
		we decide correctly $I$ w.p. at least $\frac{1}{2}$. %if and only if $w_n(S) = \OPT(U)$, which holds w.p. at least $\frac{1}{2}$. 
		It follows that $\cB$ is %by returning that $I$ has a solution if and only if $w_n(S) = B_{n,k,\alpha}$. Thus, we 
		a randomized algorithm for the $\Pi$-matroid problem which runs in polynomial time, contradicting \Cref{cor:deter}. 
		
		
		
	\end{proof}
	
	
	\subsection{Knapsack Cover with a Matroid}
	\label{sec:abs_KCM}
	We are left to prove \Cref{thm:kcm}, whose proof is similar to the previous proofs in this section.
%	\thmkcm*
	
	\begin{proof}
		Assume towards a contradiction that there is a randomized Fully PTAS $\cA$ for KCM. We use $\cA$ to decide the $\Pi$-matroid problem.  Let $I = \left(n,k,\alpha,\mm\right)$ be a $\Pi$-matroid instance. Consider the following algorithm $\cB$ that decides $I$.
		\begin{enumerate}
			\item Construct the KCM instance $%U_{n,k,\alpha}(\Pi) 
			U= ([n],\cI_{n,k,\alpha}(\Pi),w_{n},w_{n},B_{n,k,\alpha})$ along with the membership oracle of $M_{n,k,\alpha}(\Pi)$; note that the elements of $U$ are $[n] = \{1,\ldots, n \}$, the independent sets are the independent sets of the $\Pi$-matroid $\mm$ (i.e., $\cI_{n,k,\alpha}(\Pi)$), the standard weight function $w_n$ is both the cost and the size function, and the demand is simply the standard budget with parameters $n,k,\alpha$. %Let $U = U_{n,k,\alpha}(\Pi)$ for simplicity.
			\item Define $\eps = \frac{1}{6 \cdot n^3}$ and compute $\cA$ on $U$ and $\eps$; let $S$ be the returned solution.  
			\item Return that $I$ has a solution if and only if $w_n(S) = B_{n,k,\alpha}$. 
		\end{enumerate}
		
		Note that $U$ can be constructed from $I$ in time $\io$ and that $\frac{1}{\eps} = |I|^{O(1)}$; thus, the running time of $\cB$ on $I$ is $|I|^{O(1)}$, since $\cA$ is a (randomized) Fully PTAS. We now show correctness. %For the correctness, %we show that $\cB$ returns that $I$ has a solution if and only if $I$ indeed has a solution. 
		\begin{itemize}
			\item If $\cB$ returns that $I$ has a solution. Then, $w_n(S) = B_{n,k,\alpha}$; since $S$ is a solution for $U$ it also holds that $S \in \cI_{n,k,\alpha}(\Pi)$. Thus, by \Cref{cor:iff} it holds that $S$ is a solution for  $I$, and in particular $I$ has a solution.
			\item If $I$ has a solution. Then, there is $T \in \mathcal{L}_{n,k,\alpha}(\Pi)$; thus, by  \Cref{cor:iff} it follows that $w_n(T) = B_{n,k,\alpha}$ implying that $T$ is a solution for $U$.  Therefore, with probability at least $\frac{1}{2}$: 
			\begin{equation}
				\label{eq:sOPT}
				B_{n,k,\alpha} \leq \OPT(U) \leq w_n(S) \leq (1+\eps) \cdot w_n(T) \leq  w_n(T)+\frac{1}{6 \cdot n^3} \cdot 3 \cdot n^3 = B_{n,k,\alpha}+\frac{1}{2}.
			\end{equation} 	The first inequality holds since %$w_n(T) = B_{n,k,\alpha}$. In addition, $\OPT(U) \leq B_{n,k,\alpha}$ because 
			the size and the cost of each element are the same, and the size of a solution must be higher or equal from the demand. The third inequality follows (with probability at least $\frac{1}{2}$) because $S,T$ are solutions for $U$, and $S$ is computed by a randomized FPTAS with the parameter $\eps$. The fourth inequality follows from \Cref{lem:OPTu} and by the selection of $\eps$. Observe that the value of $S$ is an integer by \eqref{eq:w}; thus, by \eqref{eq:sOPT}, %w.p. at least $\frac{1}{2}$ 
			it holds that $w_n(S) = B_{n,k,\alpha}$ with probability at least $\frac{1}{2}$. Then, $\cB$ returns that $I$ has a solution with probability at least $\frac{1}{2}$.  
		\end{itemize}
		Hence, %by computing $\cA$ on $U$ and getting a solution $S$, 
		the algorithm correctly decides $I$ with probability at least $\frac{1}{2}$. %if and only if $w_n(S) = \OPT(U)$, which holds w.p. at least $\frac{1}{2}$. 
		It follows that $\cB$ is %by returning that $I$ has a solution if and only if $w_n(S) = B_{n,k,\alpha}$. Thus, we 
		a randomized algorithm for $\Pi$-matroid problem which runs in polynomial time, contradicting  \Cref{cor:deter}.

	\end{proof}
	
	
	
	
	
	
%	\subsection{Knapsack Cover with Matroid Intersection}
	